\chapter{Distributional products of three boundary currents}
\label{chap:collision-distributions}

A double pole has a distributional extension, but its Dolbeault derivative
is a contact term.  Three quadratic currents also produce a triangle of
poles.  Its three collision normals are dependent.  Partial fractions
construct an extension without multiplying those three distributions.
Complex scaling then proves that the extension preserves rational identities.

\section{The double pole and its residue}

Work on $\C$ with $z=x+iy$, $dA=dx\,dy$, and
$\partial_z=(\partial_x-i\partial_y)/2$.
Distributions act complex linearly on compactly supported smooth densities.
The residue convention is
\[
 \int^{\mathrm{res}}\alpha=\frac1{2\pi i}\int\alpha,
 \qquad d\bar z\wedge dz=2i\,dA.
\]
Fix $\kappa\in\C$, put $\lambda=-\hbar\kappa$, and give
$\hbar$ degree zero.  All formulas are polynomial in $\hbar$ on polynomial
inputs.  For $n\ge1$, define scalar distributions
\begin{equation}
 \label{eq:cd-poles}
 E_1=\frac1z\in L^1_{\mathrm{loc}}(\C),\qquad
 E_n=\frac{(-1)^{n-1}}{(n-1)!}\partial_z^{n-1}E_1,
 \qquad E_0=1.
\end{equation}
Define the $(0,1)$-currents
\[
 R_n=\frac{\pi(-1)^{n-1}}{(n-1)!}
       \partial_z^{n-1}\delta_0\,d\bar z.
\]
Here $\delta_0$ is normalized by $\langle\delta_0,\phi\,dA\rangle=\phi(0)$.

\begin{lemma}[Distributional residue]
\label{lem:cd-residue}
The restrictions of $E_n$ to $\C\setminus\{0\}$ equal $z^{-n}$.
They satisfy
\begin{equation}
 \label{eq:cd-residue}
 zE_n=E_{n-1},\qquad
 \partial_zE_n=-nE_{n+1},\qquad
 \bar\partial E_n=R_n.
\end{equation}
For every $t\in\C^\times$, pullback under $z\mapsto tz$ gives
$E_n(tz)=t^{-n}E_n(z)$.  In particular, the current contraction is
\begin{equation}
 \label{eq:cd-double-contact}
 C(z)=\lambda E_2(z),\qquad
 \bar\partial C=-\pi\lambda\partial_z\delta_0\,d\bar z.
\end{equation}
\end{lemma}
\begin{proof}
Local integrability follows from $\int_0^\epsilon r^{-1}r\,dr<\infty$.
For a test function $\phi$, apply Stokes' theorem outside $|z|=\epsilon$.
The inner circle has the clockwise orientation.  It gives
\[
 \left\langle\partial_{\bar z}E_1,\phi\,dA\right\rangle
 =-\lim_{\epsilon\downarrow0}\int_{|z|>\epsilon}
                 \frac{\partial_{\bar z}\phi}{z}\,dA
 =\pi\phi(0).
\]
The integral over the omitted disk tends to zero for the differentiated
test function.  Differentiating this identity proves the residue formula.
The derivative identity follows directly from the definition.
Also $zE_1=1$.  Differentiating this identity $n-1$ times gives
$zE_n=E_{n-1}$ for $n>1$.
Change of variables proves complex scaling for $E_1$.
Distributional differentiation gives it for every $E_n$.
\end{proof}

The coefficient $\lambda=-\hbar\kappa$ follows from the compact quantum
boundary differential.  With the conventions above, that differential is
\begin{equation}
 \label{eq:cd-boundary-carrier}
 \begin{split}
 \mathcal B_\kappa(U)&=
 \operatorname{Sym}\bigl(\Omega_c^{0,\bullet}(U)[1]\bigr)[\hbar],\\
 D&=\bar\partial+\hbar\Delta_{\mu_\kappa},\qquad
 \mu_\kappa(u,v)=-\kappa\int^{\mathrm{res}}u\wedge\partial v.
 \end{split}
\end{equation}
A function generator has degree $-1$, and a $(0,1)$-form generator has degree
zero.  The symmetric powers contain compactly supported smooth kernels
on $U^r$, with all derivative seminorms on each fixed compact support.
The algebra is a direct sum over $r$.  A point-supported current is not
an element of this smooth observable carrier.

The scalar contraction on smooth form kernels has exactly the extension
in \eqref{eq:cd-poles}.  For $u=a(w)d\bar w$ and $v=b(z)d\bar z$,
\[
 P_\kappa(u,v)=\frac{\kappa}{\pi^2}
 \left\langle E_2(z-w),a(w)b(z)\,dA_zdA_w\right\rangle.
\]
Indeed $E_2=-\partial_zE_1$, so integration by parts gives
$\kappa\pi^{-2}\int (z-w)^{-1}a(w)\partial_zb(z)\,dA_zdA_w$.
This is the Cauchy contraction in the smooth observable complex.
The Wick contraction has the opposite factor, $-\hbar P_\kappa$.
No extension of $D$ to arbitrary distributional observable kernels is used.

\section{Transverse products and the triangle}

For a scalar distribution $T$, write $\operatorname{WF}(T)$ for its real
wavefront set.  The multiplication criterion requires that no nonzero
covector of one factor is the negative of a covector of the other.
Brouder--Dang--H\'elein prove the corresponding continuity statement in
\cite[Theorem 6.1, p.~21]{BDH2014}.

For each $n\ge1$,
\begin{equation}
 \label{eq:cd-wavefront}
 \operatorname{WF}(E_n)=T_0^*\C\setminus\{0\}.
\end{equation}
The distribution is smooth away from zero and singular at zero.
Its singularity follows from its unbounded restriction $z^{-n}$.
Its wavefront set over zero is nonempty and invariant under all rotations,
by complex scaling.  Rotations act transitively on real covector directions,
which proves \eqref{eq:cd-wavefront}.

If $u,v$ are independent complex linear coordinates on $\C^2$, define
$E_a(u)E_b(v)$ as the pullback of the tensor product $E_a\otimes E_b$.
This is also the wavefront product.  Its two conormal spaces meet only at
zero.  No extension choice at $(u,v)=(0,0)$ is needed for this product.
The same construction permits one or both $E$ factors to be replaced by $R$.

Put $x=z_1-z_2$, $y=z_2-z_3$, and $c=z_3$.
Thus $z_1-z_3=x+y$.
The first shared-edge scalar coefficient is
\[
 \langle :J^2:(z_1)J(z_2)J(z_3)\rangle
       =2\lambda^2 E_2(x)E_2(x+y).
\]
The two linear coordinates $x,x+y$ are independent.
For three quadratic insertions, the connected coefficient away from the
collision divisors is
\[
 \frac{8\lambda^3}{x^2y^2(x+y)^2}.
\]
At $x=y=0$, nonzero covectors from the three conormal spaces sum to zero:
choose $\xi\,dx$, $\xi\,dy$, and $-\xi\,d(x+y)$ in real covector notation.
Thus the wavefront multiplication criterion fails for the three factors.
Failure of this sufficient criterion alone does not prove that every
regularization fails.  The following formula specifies the extension.

\begin{proposition}[Triangle extension]
\label{prop:cd-triangle}
The distribution
\begin{equation}
 \label{eq:cd-triangle}
 \begin{split}
 T_\triangle={}&E_2(x)E_4(y)+E_2(x+y)E_4(y)\\
 &-2E_1(x)E_5(y)+2E_1(x+y)E_5(y)
 \end{split}
\end{equation}
extends $x^{-2}y^{-2}(x+y)^{-2}$.  Each product displayed is transverse.
It has complex scaling degree $-6$:
$T_\triangle(tx,ty)=t^{-6}T_\triangle(x,y)$.
It is independent of the chosen partial-fraction decomposition, provided
that each smooth point of a collision divisor uses its Laurent extension
\eqref{eq:cd-poles} and each term has this exact complex scaling.
\end{proposition}
\begin{proof}
Ordinary partial fractions give
\[
 \frac1{x^2y^2(x+y)^2}
 =\frac1{x^2y^4}+\frac1{(x+y)^2y^4}
       -\frac2{xy^5}+\frac2{(x+y)y^5}.
\]
Every pair of denominators on the right is linearly independent.
The preceding tensor construction therefore defines the right side
as a distribution.  The formulas agree on the complement of the divisors.
They also agree near every smooth point of a divisor under the stated
Laurent prescription.  A difference between two such extensions is supported
at the origin.  Both satisfy the displayed complex scaling.
Lemma~\ref{lem:cd-scaling-uniqueness} below makes their difference zero.
\end{proof}

\begin{lemma}[Complex scaling excludes supported ambiguity]
\label{lem:cd-scaling-uniqueness}
Let $d\ge1$ and $m\in\Z$.
A distribution $S$ on $\C^d$, supported at zero, cannot satisfy
$S(tz)=t^{-m}S(z)$ for every $t\in\C^\times$, unless $S=0$.
\end{lemma}
\begin{proof}
A distribution supported at zero is a finite sum
$\sum_{\alpha,\beta}s_{\alpha\beta}
\partial_z^\alpha\partial_{\bar z}^\beta\delta_0$.
To obtain this description, use its finite order on a compact neighborhood.
Taylor subtraction shows that it vanishes on test functions with sufficiently
many zero derivatives at zero.  It therefore factors through that finite
jet space, whose dual consists of the displayed derivatives.

Write $a=|\alpha|$ and $b=|\beta|$.
Under $t=r>0$, a displayed derivative scales by $r^{-2d-a-b}$.
Under $t=e^{i\theta}$, it scales by $e^{-i(a-b)\theta}$.
Linear independence of these radial and angular characters requires
\[
 2d+a+b=m,\qquad a-b=m
\]
for every nonzero coefficient.  Subtraction gives $2d+2b=0$, which is
impossible.  Hence every coefficient vanishes.
\end{proof}

Real scaling alone would permit second derivatives of $\delta_0$ in the
triangle example.  Their angular weights fail the required weight $-6$.
The exact complex scaling is therefore a substantive extension condition.
The general scaling-degree extension theorem and its supported ambiguity
are given by Brunetti--Fredenhagen \cite[Theorems 5.2--5.3, pp.~23--24]{BF1999}.
The proof above also excludes ambiguities at higher pole orders in this
holomorphic setting.

\section{All rational coefficients at three insertions}

Let
\[
 R=\C[x,y,x^{-1},y^{-1},(x+y)^{-1}],\qquad
 R_3=\C[c]\otimes R.
\]
Grade $R$ by simultaneous complex dilation of $x,y$.
Every element is a finite sum of homogeneous rational functions.

\begin{theorem}[Three-insertion distributional extension]
\label{thm:cd-three-extension}
There is a unique complex linear injection
\[
 \mathcal E_3:R_3\longrightarrow\mathcal D'(\C^3)
\]
with these properties.  It extends each rational function away from its
collision divisors.  It is polynomial linear in $c$.
Near a smooth point of a divisor it uses the transverse Laurent extension
$z^{-n}\mapsto E_n(z)$.  On each homogeneous relative-coordinate component
of degree $-m$, it has exact complex scaling $t^{-m}$.

It commutes with holomorphic derivatives, multiplication by holomorphic
polynomials, simultaneous translations, and permutations of the insertions.
On products with independent collision normals it agrees with the tensor
product construction.  Formula~\eqref{eq:cd-triangle} is its value on the
triangle coefficient.
\end{theorem}
\begin{proof}
It is enough to construct the extension of a homogeneous summand.
For positive integers $a,b$ and $c_0\ge0$, the identity
\[
 \frac1{x^ay^b(x+y)^{c_0}}
 =\frac1{x^{a-1}y^b(x+y)^{c_0+1}}
  +\frac1{x^ay^{b-1}(x+y)^{c_0+1}}
\]
reduces $a+b$ by one in each term.  Repetition expresses the denominator
as a finite sum involving at most two independent linear factors.
Polynomial numerators can be multiplied into these distributions.
This constructs a finite sum of pullbacks of $E_p\otimes E_q$,
possibly with smooth polynomial factors.  The homogeneity is unchanged.

The construction has the prescribed Laurent extension away from the
origin.  To check this assertion, use the singular linear form as a local
coordinate $u$.  Every other denominator is holomorphic and invertible
there.  For a holomorphic multiplier $h$, divide its Taylor expansion as
$h(u)=\sum_{j=0}^{n-1}h_j u^j+u^n h_n(u)$.
The identities $u^jE_n(u)=E_{n-j}(u)$ for $j<n$, and
$u^nE_n(u)=1$, give exactly its Laurent expansion as a distribution.
The coefficients can depend holomorphically on the other coordinate.

Two partial-fraction decompositions therefore differ by a distribution
supported at $(x,y)=0$.  They have the same complex homogeneity.
Lemma~\ref{lem:cd-scaling-uniqueness}, with $d=2$, proves equality.
This also proves linearity within each homogeneous component and independence
of polynomial numerator presentations.  Sum the components and then extend
polynomial linearly in $c$.  Restriction to the dense complement of the
divisors proves injectivity.

A holomorphic derivative of the constructed extension has the Laurent
prescription of the differentiated rational function.  Its homogeneity
is the differentiated degree.  The same uniqueness argument proves
commutation with relative holomorphic derivatives.  The argument also
applies to multiplication by homogeneous polynomials.  Linearity gives
all polynomial multipliers and all holomorphic derivatives, including $c$.

For a rational function of differences, every permutation acts by an
invertible complex linear map of $(x,y)$ preserving the three divisors.
Pullback preserves both the Laurent prescription and complex homogeneity,
so uniqueness proves covariance.  Polynomial multiplication then gives
covariance on $R_3$, since a permuted $c$ equals $c$ plus a linear
combination of $x,y$.  Simultaneous translation changes only $c$.
The transverse tensor prescription has the same defining properties,
which proves the remaining assertion.
\end{proof}

The image $\mathscr R_3=\mathcal E_3(R_3)$ carries the specified product
\begin{equation}
 \label{eq:cd-renormalized-product}
 \mathcal E_3(f)\mathbin{\circ}\mathcal E_3(g)=\mathcal E_3(fg).
\end{equation}
Injectivity makes this a well-defined associative algebra over $\C$.
This product is a prescription on $\mathscr R_3$.  It does not assert
multiplication of arbitrary distributions.  In particular, adjoining the
residue currents requires further product definitions.

If $\Gamma$ is the union of the conormal bundles of the three pair diagonals
and the full conormal bundle of the total diagonal, then
\[
 \operatorname{WF}(\mathcal E_3(f))\subseteq\Gamma\setminus\{0\}.
\]
The bound follows termwise from the transverse tensor formula.
At a point where only one divisor is present, only its normal covectors
occur.  At the total diagonal all relative covectors are allowed.
No covector in the simultaneous translation direction occurs.

\section{Contact terms and three-current compatibility}

Pull back $R_n$ with its displayed $d\bar z$ factor under each linear
coordinate map.  Distributional differentiation gives the explicit triangle
contact current
\begin{equation}
 \label{eq:cd-triangle-contact}
 \begin{split}
 \bar\partial T_\triangle={}&R_2(x)E_4(y)+E_2(x)R_4(y)\\
 &+R_2(x+y)E_4(y)+E_2(x+y)R_4(y)\\
 &-2R_1(x)E_5(y)-2E_1(x)R_5(y)\\
 &+2R_1(x+y)E_5(y)+2E_1(x+y)R_5(y).
 \end{split}
\end{equation}
Every term is a transverse tensor product.  Thus the equation includes
the total collision and does not require multiplication on one divisor.
For independent $u,v$,
\[
 \bar\partial\bigl(R_a(u)E_b(v)+E_a(u)R_b(v)\bigr)
 =-R_a(u)\wedge R_b(v)+R_a(u)\wedge R_b(v)=0.
\]
Consequently \eqref{eq:cd-triangle-contact} is closed, including its contact
terms on the total diagonal.

To obtain actual compact observable outputs, choose
$0<r<\rho<R$, put $A=D_r$, $U=D_R$, and choose a smooth radial cutoff $\chi$.
Require $\chi=1$ on $D_\rho$ and $\operatorname{supp}\chi\Subset D_R$.
For $a\in A$, define
\[
 e_a(z)=-\frac{\bar\partial_z\chi(z)}{z-a},\qquad
 e_{a,n}(z)=\partial_a^n e_a(z).
\]
These are compactly supported smooth $(0,1)$-forms on a fixed outer annulus.
They depend holomorphically on $a$ in every smooth seminorm.
No support shrinks to a point as $a$ moves.

\begin{proposition}[Compactly smeared collision products]
\label{prop:cd-smeared}
For three elementary insertions, the expression
\begin{equation}
 \label{eq:cd-three-current}
 \mathscr W_3(a,b,c)=e_ae_be_c+
 \lambda\bigl(E_2(a-b)e_c+E_2(a-c)e_b+E_2(b-c)e_a\bigr)
\end{equation}
defines a continuous linear map
\[
 C_c^\infty(A^3,\operatorname{Dens})[\hbar]
       \longrightarrow\mathcal B_\kappa(U)^0,
 \qquad F\longmapsto\langle\mathscr W_3,F\rangle.
\]
Its values are $D$-closed compact smooth observable kernels.
More generally, the three-insertion Wick formula for any finite monomials
in the jets $e_{a,n}$ defines such a map by applying $\mathcal E_3$
to each rational contraction coefficient.
The field number of the output is finite.

For the elementary product, the external Dolbeault derivative is
\begin{equation}
 \label{eq:cd-three-Ward}
 \bar\partial_{a,b,c}\mathscr W_3
 =\lambda\bigl(R_2(a-b)e_c+R_2(a-c)e_b+R_2(b-c)e_a\bigr).
\end{equation}
Here, for example,
$R_2(a-b)=-\pi\partial_{a-b}\delta(a-b)(d\bar a-d\bar b)$.
The two parenthesizations of the complete three-insertion Wick product
have zero difference as distributions under $\mathcal E_3$.
\end{proposition}
\begin{proof}
A product of the $e$ forms is a smooth kernel on a fixed product of
compact annuli.  For each compact parameter support, all parameter and
field derivatives are bounded in every required seminorm.
Each coefficient distribution has finite order on this support.
Applying it to the parameter variables therefore gives a smooth kernel
in the remaining field variables, with the same annular support.
For example, a bound of distribution order $N$ gives
\[
 \|\langle T,F e\rangle\|_{C^k}
 \le C_{T,K,k}\|F\|_{C^N}
       \sup_K\|e\|_{C^{N+k}}.
\]
The same estimate applies to finitely many field factors.
It proves continuity on every fixed compact support, hence on the test
function direct limit.  This construction also handles the holomorphic
smooth coefficients multiplying the rational Wick coefficients.
There are finitely many matchings, so the sum has finite field number.

Every output factor has antiholomorphic degree one.
Its internal Dolbeault derivative is zero, and the quantum cocycle vanishes
on pairs of such factors.  Thus $D$ kills the output smooth kernels.
This uses the fixed differential \eqref{eq:cd-boundary-carrier}.
Holomorphic dependence of the $e$ forms and
Lemma~\ref{lem:cd-residue} give \eqref{eq:cd-three-Ward}.

For any polynomial inputs, label each individual current factor.
A Wick matching pairs factors in different insertion groups.
Grouping first the first two groups, or first the last two, enumerates
the same matchings with the same coefficients.  Therefore the complete
rational expressions agree before extension.  Linearity of $\mathcal E_3$
preserves this equality.  Multiplication by the same smooth unpaired
kernels and smearing preserve it as well.
\end{proof}

For separated centers, choose disjoint inner cutoff disks and set
$u_a=-\bar\partial\chi_a/(z-a)$.
The compact quantum homotopy identifies the product $u_au_bu_c$ with
\eqref{eq:cd-three-current} in cohomology.
This gives the separated boundary comparison with the fixed sign
$\lambda=-\hbar\kappa$.
At collisions, Proposition~\ref{prop:cd-smeared} extends its smooth annular
output as an insertion-parameter distribution.  It does not identify a
Dirac current with a smooth observable generator.
Equation~\eqref{eq:cd-three-Ward} specifies the contact obstruction to
regarding this extension alone as a Dolbeault chain map.

The regular vertex product has a different three-point identity.
With $q_{n+1}=[-\bar\partial\chi/z^{n+1}]$, the field on polynomial states is
\[
 J(z)=\sum_{m\ge1}q_mz^{m-1}
       +\lambda\sum_{m\ge1}m\partial_{q_m}z^{-m-1}.
\]
The regular product $a_{(-1)}b$ is the coefficient of $z^0$ in $Y(a,z)b$.
Direct normal ordering gives
\[
 \begin{split}
 q_{1(-1)}q_1&=q_1^2,\qquad q_{1(-1)}q_1^2=q_1^3,\\
 Y(q_1^2,z)q_1&=J_-(z)^2q_1+2\lambda z^{-2}J_-(z).
 \end{split}
\]
The constant coefficient of the last expression is $q_1^3+2\lambda q_3$.
Consequently its first associativity defect is
\begin{equation}
 \label{eq:cd-regular-associator}
 (q_{1(-1)}q_1)_{(-1)}q_1-q_{1(-1)}(q_{1(-1)}q_1)
 =2\lambda q_3=-2\hbar\kappa q_3.
\end{equation}
A rational coefficient extension on the insertion space does not identify
iterated Laurent constant terms with multiplication at a point.
Thus \eqref{eq:cd-regular-associator} is compatible with the zero grouping
defect in Proposition~\ref{prop:cd-smeared}.

\section{Smooth interval coefficients and the remaining chain condition}

Let $I=[0,1]$ have its increasing orientation.
Use the full smooth de Rham complex
$\Omega^\bullet(I)=C^\infty(I)\oplus C^\infty(I)\,dt$,
where smooth functions extend to a neighborhood of the closed interval.
There are no endpoint vanishing conditions.
The finite rational Wick coefficients in three insertion variables form
$R_3[x_{i,n}][\hbar]$, with each polynomial involving finitely many jets.
The distributional coefficient map is
\[
 1\otimes\mathcal E_3:
 \Omega^\bullet(I)\otimes R_3[x_{i,n}][\hbar]
 \longrightarrow
 \Omega^\bullet(I)\otimes\mathscr R_3[x_{i,n}][\hbar].
\]
For jointly smooth interval dependence, apply $\mathcal E_3$ to each
fixed rational coefficient and retain its full smooth interval amplitude.
The finite-order estimate in Proposition~\ref{prop:cd-smeared} permits
interval derivatives, endpoint evaluation, and integration to pass through
the parameter distribution.  The map therefore commutes with $d_I$.
It preserves every three-insertion identity formed from interval-linear
maps and the prescribed rational Wick product.
This statement includes all smooth interval amplitudes, independently
of any polynomial subcomplex.

All three-insertion formulas extend coefficientwise to $\hbar$-adic series
whose individual coefficients have finite polynomial degree and finitely
many jets.  A fixed output coefficient receives finitely many input
coefficients and matching terms.  No uniform distribution order or field
number across the entire formal series is asserted.

If a current convention uses $+\hbar\kappa_I(z-w)^{-2}$, the boundary
comparison requires $\kappa_I=-\kappa$.
All jet contractions then become
$\lambda\partial_{z_i}^r\partial_{z_j}^sE_2(z_i-z_j)$ with the same
undivided-jet convention.  Divided jets add the factors $(r!s!)^{-1}$.

On insertion-parameter forms with observable values, the total differential is
\[
 \mathbb D(\omega\otimes b)
 =\bar\partial\omega\otimes b+(-1)^{|\omega|}\omega\otimes Db.
\]
It squares to zero.  Formula~\eqref{eq:cd-three-Ward} is the value of
$\mathbb D$ on the elementary extended product.
The corresponding triangle contact is \eqref{eq:cd-triangle-contact}.
To construct operations on collision chains, one must supply maps for
these residue currents and prove compatibility with their boundary strata.
The associative product \eqref{eq:cd-renormalized-product} does not itself
supply products involving these $(0,1)$-currents.
Nor does the three-insertion construction specify extensions for four or
more insertion variables.  These are separate mathematical obligations.
